% =====================================================================
%  MediSync — Présentation Frontend Angular
%  Document de soutenance — Génie Informatique S4
%  Développement Web & Mobile - Projet P5
% =====================================================================
\documentclass[11pt,a4paper]{article}

% ---------- Encodage & langue ----------
\usepackage[utf8]{inputenc}
\usepackage[T1]{fontenc}
\usepackage[french]{babel}
\usepackage{lmodern}

% ---------- Mise en page ----------
\usepackage[margin=2.2cm]{geometry}
\usepackage{setspace}
\onehalfspacing
\usepackage{parskip}
\usepackage{microtype}

% ---------- Couleurs & Boîtes ----------
\usepackage{xcolor}
\definecolor{primary}{HTML}{0D6EFD}
\definecolor{secondary}{HTML}{06B6D4}
\definecolor{accent}{HTML}{8B5CF6}
\definecolor{darkgray}{HTML}{0F172A}
\definecolor{mediumgray}{HTML}{475569}
\definecolor{lightgray}{HTML}{F8FAFC}
\definecolor{bordergray}{HTML}{E2E8F0}
\definecolor{success}{HTML}{10B981}
\definecolor{warning}{HTML}{F59E0B}
\definecolor{danger}{HTML}{EF4444}

\usepackage[most]{tcolorbox}
\tcbuselibrary{skins,breakable}

% Boîte d'information
\newtcolorbox{infobox}[1][]{
  enhanced, breakable,
  colback=primary!5, colframe=primary,
  boxrule=0pt, leftrule=4pt,
  arc=4pt, outer arc=4pt,
  fonttitle=\bfseries, coltitle=primary,
  left=10pt, right=10pt, top=8pt, bottom=8pt,
  #1
}

% Boîte de succès
\newtcolorbox{successbox}[1][]{
  enhanced, breakable,
  colback=success!5, colframe=success,
  boxrule=0pt, leftrule=4pt,
  arc=4pt, outer arc=4pt,
  fonttitle=\bfseries, coltitle=success,
  left=10pt, right=10pt, top=8pt, bottom=8pt,
  #1
}

% Boîte de feature
\newtcolorbox{featurebox}[1][]{
  enhanced, breakable,
  colback=white, colframe=bordergray,
  boxrule=1pt, arc=8pt, outer arc=8pt,
  left=14pt, right=14pt, top=10pt, bottom=10pt,
  #1
}

% ---------- Titres ----------
\usepackage{titlesec}
\titleformat{\section}
  {\Large\bfseries\color{primary}}
  {\thesection}{1em}{}
  [\vspace{2pt}{\color{primary}\rule{\textwidth}{1.5pt}}]
\titleformat{\subsection}
  {\large\bfseries\color{darkgray}}
  {\thesubsection}{0.8em}{}
\titleformat{\subsubsection}
  {\normalsize\bfseries\color{mediumgray}}
  {\thesubsubsection}{0.6em}{}

% ---------- Listings code ----------
\usepackage{listings}
\lstdefinelanguage{TypeScript}{
  keywords={typeof, new, true, false, catch, function, return, null, switch,
            var, if, in, while, do, else, case, break, let, const, async,
            await, import, from, export, default, class, extends, implements,
            interface, public, private, readonly, static, this, signal,
            computed, inject, Injectable, Component, Input, Output},
  keywordstyle=\color{primary}\bfseries,
  ndkeywords={String, number, boolean, void, any, Array, Promise, Observable,
              Routes, CanActivateFn, HttpInterceptorFn, User, UserRole},
  ndkeywordstyle=\color{accent}\bfseries,
  identifierstyle=\color{darkgray},
  sensitive=true,
  comment=[l]{//},
  morecomment=[s]{/*}{*/},
  commentstyle=\color{mediumgray}\itshape,
  stringstyle=\color{success},
  morestring=[b]',
  morestring=[b]"
}

\lstdefinestyle{codestyle}{
  language=TypeScript,
  backgroundcolor=\color{lightgray},
  basicstyle=\ttfamily\footnotesize\color{darkgray},
  frame=single,
  rulecolor=\color{bordergray},
  framerule=1pt,
  framesep=8pt,
  xleftmargin=10pt,
  xrightmargin=10pt,
  showstringspaces=false,
  breaklines=true,
  numbers=left,
  numberstyle=\tiny\color{mediumgray},
  numbersep=8pt,
  tabsize=2,
  captionpos=b
}
\lstset{style=codestyle}

% ---------- Tableaux & listes ----------
\usepackage{array}
\usepackage{booktabs}
\usepackage{tabularx}
\usepackage{colortbl}
\usepackage{enumitem}
\setlist[itemize]{leftmargin=*,itemsep=4pt}

% ---------- Liens ----------
\usepackage{hyperref}
\hypersetup{
  colorlinks=true,
  linkcolor=primary,
  urlcolor=primary,
  citecolor=accent,
  pdftitle={MediSync - Frontend Angular},
  pdfauthor={Equipe MediSync}
}

% ---------- Graphiques ----------
\usepackage{graphicx}
\usepackage{float}
\usepackage{tikz}
\usetikzlibrary{shapes,arrows,positioning,shadows,fit,backgrounds}

% =====================================================================
%  DEBUT DU DOCUMENT
% =====================================================================
\begin{document}

% ============================================
%  PAGE DE GARDE
% ============================================
\begin{titlepage}
\begin{tikzpicture}[remember picture, overlay]
  \fill[primary] (current page.north west) rectangle ([yshift=-3cm]current page.north east);
  \fill[secondary] ([yshift=-3cm]current page.north west) rectangle ([yshift=-3.3cm]current page.north east);
\end{tikzpicture}

\vspace{0.5cm}
\begin{center}
{\color{white}\Huge\bfseries MediSync}\\[0.3cm]
{\color{white}\Large Système de Gestion de Clinique Médicale}
\end{center}

\vspace{3cm}

\begin{center}
\begin{tcolorbox}[
  enhanced, colback=white, colframe=primary,
  boxrule=2pt, arc=12pt, width=0.85\textwidth,
  drop shadow={shadow xshift=0pt, shadow yshift=-4pt, opacity=0.15}
]
\centering
{\color{primary}\Huge\bfseries Partie Frontend}\\[0.4cm]
{\Large\color{darkgray} Architecture, Conception \& Réalisation}\\[0.6cm]
{\large\color{mediumgray} avec Angular 21 \& Angular Material}\\[1cm]

{\color{primary}\rule{0.5\textwidth}{1pt}}\\[0.5cm]

{\large\bfseries Soutenance de Projet}\\[0.2cm]
{\normalsize Développement Web \& Mobile - P5}\\[0.2cm]
{\normalsize Génie Informatique - Semestre 4}\\[0.6cm]

{\color{mediumgray}\large 03 Juin 2026}
\end{tcolorbox}
\end{center}

\vfill

\begin{tikzpicture}[remember picture, overlay]
  \fill[primary] ([yshift=2.5cm]current page.south west) rectangle (current page.south east);
\end{tikzpicture}

\begin{center}
\vspace{1cm}
{\color{white}\large\bfseries Équipe MediSync}\\[0.2cm]
{\color{white} Génie Informatique S4 \textbar{} Année universitaire 2025-2026}
\end{center}

\end{titlepage}

% ============================================
%  TABLE DES MATIERES
% ============================================
\renewcommand{\contentsname}{\color{primary}\textbf{Sommaire}}
\tableofcontents
\thispagestyle{empty}
\newpage

% ============================================
%  1. INTRODUCTION
% ============================================
\section{Introduction}

\subsection{Présentation du projet MediSync}

\textbf{MediSync} est une plateforme web et mobile dédiée à la digitalisation
de la gestion d'une clinique médicale multi-praticiens. Le système centralise
la prise de rendez-vous, la gestion des dossiers patients, les ordonnances
électroniques et le pilotage administratif au sein d'une expérience utilisateur
cohérente et sécurisée.

\begin{infobox}[title={Objectifs métiers}]
\begin{itemize}
  \item \textbf{Patients} : prise de rendez-vous, accès au dossier médical, suivi des ordonnances.
  \item \textbf{Personnel médical} : gestion du planning, comptes rendus, prescriptions électroniques.
  \item \textbf{Secrétariat} : création de comptes patients, facturation.
  \item \textbf{Administration} : pilotage global, statistiques, conformité.
\end{itemize}
\end{infobox}

\subsection{Périmètre de la partie frontend}

La partie frontend constitue \textbf{l'interface utilisateur principale}
du système. Elle est responsable de :

\begin{itemize}
  \item Présenter l'application au grand public via une \textbf{landing page} engageante.
  \item Assurer l'\textbf{authentification} et la gestion des sessions utilisateurs.
  \item Fournir un \textbf{tableau de bord adapté} à chacun des quatre rôles métier.
  \item Communiquer avec l'API REST sécurisée (\textbf{Spring Boot}) via JWT.
  \item Garantir une expérience utilisateur \textbf{responsive}, accessible et fluide.
\end{itemize}

\newpage

% ============================================
%  2. STACK TECHNIQUE
% ============================================
\section{Stack technique}

\subsection{Choix technologiques}

\begin{table}[H]
\centering
\renewcommand{\arraystretch}{1.5}
\begin{tabularx}{\textwidth}{>{\columncolor{lightgray}}lX}
\toprule
\rowcolor{primary!10}
\textbf{Couche} & \textbf{Technologie retenue} \\
\midrule
Framework & \textbf{Angular 21} (composants standalone, signaux) \\
Langage & \textbf{TypeScript 5.7} avec mode strict \\
UI Components & \textbf{Angular Material 21} (Material Design 3, thème \emph{azure-blue}) \\
Styles & \textbf{SCSS} avec variables CSS et design tokens \\
Routing & \textbf{Angular Router} avec lazy-loading et route guards \\
Formulaires & \textbf{Reactive Forms} avec validateurs personnalisés \\
HTTP & \texttt{HttpClient} avec intercepteurs JWT \\
Charts & \textbf{Chart.js + ng2-charts} \\
Animations & CSS keyframes + IntersectionObserver \\
\bottomrule
\end{tabularx}
\caption{Stack technique du frontend}
\end{table}

\subsection{Pourquoi Angular ?}

Angular a été \textbf{imposé par le cahier des charges}, mais son adoption se
justifie pleinement :

\begin{itemize}
  \item \textbf{Framework opinionated} : structure claire pour un projet d'équipe.
  \item \textbf{TypeScript natif} : sécurité de types essentielle dans un contexte médical.
  \item \textbf{Signaux et standalone components} : performances optimales (Angular 16+).
  \item \textbf{Écosystème mature} : Angular Material couvre 95\% des besoins UI.
  \item \textbf{RxJS} : gestion réactive des flux asynchrones (API, formulaires).
\end{itemize}

\newpage

% ============================================
%  3. ARCHITECTURE
% ============================================
\section{Architecture frontend}

\subsection{Organisation modulaire}

L'application est organisée selon le pattern \textbf{Core / Shared / Features},
recommandé par la communauté Angular pour les projets de taille moyenne à
grande.

\begin{lstlisting}[language={}, basicstyle=\ttfamily\footnotesize]
src/app/
|-- core/                  <- Singletons & services applicatifs
|   |-- auth/
|   |-- guards/            <- authGuard, roleGuard, guestGuard
|   |-- interceptors/      <- jwtInterceptor
|   |-- services/          <- AuthService
|   |-- models/            <- User, AuthResponse, UserRole
|
|-- shared/                <- Composants & directives reutilisables
|   |-- components/        <- PageHeader
|   |-- directives/        <- RevealOnScrollDirective
|
|-- features/              <- Decoupage par domaine metier
|   |-- landing/           <- Page d'accueil publique
|   |-- auth/              <- Login, Register
|   |-- patient/           <- Dashboard, RDV, Dossier, Ordonnances, Profil
|   |-- doctor/            <- Dashboard, Planning, Consultations, Patients
|   |-- secretary/         <- Dashboard, RDV, Facturation
|   |-- admin/             <- Dashboard, Etablissement, Personnel,
|                              Finance, Statistiques
|
|-- layouts/               <- Layouts d'application
|   |-- main-layout/       <- Shell avec sidenav + toolbar
|
|-- app.config.ts          <- Providers (router, http, animations)
|-- app.routes.ts          <- Routes lazy-loadees
\end{lstlisting}

\begin{successbox}[title={Bénéfices de cette structure}]
\begin{itemize}
  \item \textbf{Séparation claire des responsabilités} entre logique transverse et features métier.
  \item \textbf{Lazy-loading par feature} : chaque rôle charge uniquement son code.
  \item \textbf{Évolutivité} : ajout d'un nouveau module sans toucher l'existant.
  \item \textbf{Testabilité} : composants isolés et services injectables.
\end{itemize}
\end{successbox}

\subsection{Diagramme d'architecture}

\begin{figure}[H]
\centering
\begin{tikzpicture}[
  node distance=0.6cm,
  block/.style={rectangle, draw=primary, fill=primary!10, rounded corners=4pt,
                minimum width=3.4cm, minimum height=0.9cm, font=\small\bfseries,
                text=darkgray, drop shadow={shadow xshift=0pt, shadow yshift=-1pt, opacity=0.15}},
  service/.style={rectangle, draw=accent, fill=accent!10, rounded corners=4pt,
                  minimum width=3.4cm, minimum height=0.9cm, font=\small\bfseries,
                  text=darkgray},
  guard/.style={rectangle, draw=warning, fill=warning!10, rounded corners=4pt,
                minimum width=3.4cm, minimum height=0.7cm, font=\small,
                text=darkgray},
  backend/.style={rectangle, draw=success, fill=success!10, rounded corners=4pt,
                  minimum width=4cm, minimum height=1cm, font=\small\bfseries,
                  text=darkgray},
  arrow/.style={->, thick, draw=mediumgray}
]

\node[block] (landing) {Landing Page};
\node[block, right=of landing] (auth) {Login / Register};
\node[block, right=of auth] (dashboard) {Dashboards (4 rôles)};

\node[guard, below=of landing] (guestG) {guestGuard};
\node[guard, below=of auth] (authG) {authGuard};
\node[guard, below=of dashboard] (roleG) {roleGuard};

\node[service, below=1cm of authG] (authS) {AuthService (signals)};
\node[service, right=of authS] (jwtI) {jwtInterceptor};

\node[backend, below=1cm of authS] (api) {API REST Spring Boot};

\draw[arrow] (landing) -- (guestG);
\draw[arrow] (auth) -- (authG);
\draw[arrow] (dashboard) -- (roleG);
\draw[arrow] (guestG) -- (authS);
\draw[arrow] (authG) -- (authS);
\draw[arrow] (roleG) -- (authS);
\draw[arrow] (authS) -- (jwtI);
\draw[arrow] (jwtI) -- (api);
\draw[arrow] (authS) -- (api);

\end{tikzpicture}
\caption{Flux d'authentification et de protection des routes}
\end{figure}

\newpage

% ============================================
%  4. LANDING PAGE
% ============================================
\section{Landing Page}

\subsection{Objectif et conception}

La landing page est le \textbf{premier point de contact} avec les utilisateurs.
Elle doit transmettre en quelques secondes la valeur de MediSync et inciter
à l'action (inscription ou connexion).

\subsection{Sections}

\begin{enumerate}
  \item \textbf{Hero} avec titre accrocheur, illustration animée et 2 CTA (Créer un compte / Se connecter).
  \item \textbf{Stats} : 4 indicateurs clés (50K+ patients, 320+ praticiens, 99.9\% uptime, 24/7 support).
  \item \textbf{Fonctionnalités} : 6 cartes représentant les modules principaux.
  \item \textbf{How it works} : 3 étapes pour réserver un rendez-vous.
  \item \textbf{Témoignages} : 3 avis (médecin, patient, directrice).
  \item \textbf{Call-to-action final} avec dégradé immersif.
  \item \textbf{Footer} avec liens utiles et réseaux sociaux.
\end{enumerate}

\subsection{Animations \& interactivité}

\begin{featurebox}
\textbf{Animations CSS pures} (aucune librairie externe) :
\begin{itemize}
  \item \textbf{Blobs gradients} flottants en arrière-plan du hero (\emph{keyframes blob}).
  \item \textbf{Pulse} sur le badge "Certifié RGPD" (\emph{double pulse-ring}).
  \item \textbf{Gradient text} animé sur le mot clé du titre.
  \item \textbf{Floating cards} survolant l'image principale.
  \item \textbf{Hover lifts} sur les cartes feature/testimonial.
\end{itemize}
\end{featurebox}

\begin{featurebox}
\textbf{Scroll-reveal via une directive personnalisée} :
\begin{lstlisting}[language=TypeScript]
@Directive({ selector: '[appReveal]', standalone: true })
export class RevealOnScrollDirective implements OnInit {
  @Input() revealDelay = 0;

  ngOnInit() {
    const node = this.el.nativeElement;
    node.classList.add('reveal');
    this.observer = new IntersectionObserver((entries) => {
      entries.forEach((e) => {
        if (e.isIntersecting) {
          e.target.classList.add('revealed');
          this.observer?.unobserve(e.target);
        }
      });
    }, { threshold: 0.15 });
    this.observer.observe(node);
  }
}
\end{lstlisting}
\end{featurebox}

\newpage

% ============================================
%  5. AUTHENTIFICATION
% ============================================
\section{Système d'authentification}

\subsection{Composants}

\begin{table}[H]
\centering
\renewcommand{\arraystretch}{1.4}
\begin{tabularx}{\textwidth}{lX}
\toprule
\rowcolor{primary!10}
\textbf{Élément} & \textbf{Rôle} \\
\midrule
\texttt{LoginComponent} & Formulaire de connexion avec validation reactive \\
\texttt{RegisterComponent} & Création de compte avec règles de mot de passe fort \\
\texttt{AuthService} & Service singleton avec état réactif (signals) \\
\texttt{jwtInterceptor} & Injection automatique du Bearer token \\
\texttt{authGuard} & Protection des routes authentifiées \\
\texttt{roleGuard} & Vérification du rôle utilisateur \\
\texttt{guestGuard} & Empêche les utilisateurs connectés d'accéder à /auth \\
\bottomrule
\end{tabularx}
\caption{Composants du système d'authentification}
\end{table}

\subsection{AuthService basé sur les signals}

\begin{lstlisting}[language=TypeScript, caption={Extrait d'AuthService}]
@Injectable({ providedIn: 'root' })
export class AuthService {
  private readonly http = inject(HttpClient);
  private readonly router = inject(Router);

  private readonly userSignal = signal<User | null>(this.loadUser());

  readonly user = this.userSignal.asReadonly();
  readonly isAuthenticated = computed(() => this.userSignal() !== null);
  readonly role = computed(() => this.userSignal()?.role ?? null);

  login(payload: LoginRequest): Observable<AuthResponse> {
    return this.http
      .post<AuthResponse>(`${environment.apiUrl}/auth/login`, payload)
      .pipe(tap((res) => this.persist(this.normalize(res))));
  }

  logout(): void {
    localStorage.removeItem(TOKEN_KEY);
    localStorage.removeItem(USER_KEY);
    this.userSignal.set(null);
    this.router.navigate(['/auth/login']);
  }
}
\end{lstlisting}

\subsection{Validateur de mot de passe fort}

Conforme au cahier des charges (article 3.1.1) :
\textbf{8 caractères minimum}, \textbf{1 majuscule}, \textbf{1 chiffre},
\textbf{1 caractère spécial}.

\begin{lstlisting}[language=TypeScript]
function strongPassword(control: AbstractControl): ValidationErrors | null {
  const v = control.value as string;
  if (!v) return null;
  const ok = v.length >= 8 &&
             /[A-Z]/.test(v) &&
             /[0-9]/.test(v) &&
             /[^A-Za-z0-9]/.test(v);
  return ok ? null : { weak: true };
}
\end{lstlisting}

\newpage

% ============================================
%  6. ROUTING
% ============================================
\section{Routing et protection des accès}

\subsection{Stratégie}

\begin{itemize}
  \item \textbf{Lazy-loading systématique} via \texttt{loadComponent}.
  \item \textbf{Layout partagé} (\texttt{MainLayoutComponent}) injecté en
        composant parent pour les routes authentifiées.
  \item \textbf{Guards composables} : authGuard + roleGuard.
\end{itemize}

\subsection{Routes principales}

\begin{table}[H]
\centering
\renewcommand{\arraystretch}{1.4}
\begin{tabularx}{\textwidth}{>{\ttfamily}lX}
\toprule
\rowcolor{primary!10}
\textbf{URL} & \textbf{Protection / Description} \\
\midrule
/ & Public — landing page \\
/auth/login & guestGuard \\
/auth/register & guestGuard \\
/patient/* & authGuard + roleGuard(['patient']) \\
/doctor/* & authGuard + roleGuard(['doctor']) \\
/secretary/* & authGuard + roleGuard(['secretary']) \\
/admin/* & authGuard + roleGuard(['admin']) \\
\bottomrule
\end{tabularx}
\caption{Cartographie des routes}
\end{table}

\subsection{Extrait de configuration des routes}

\begin{lstlisting}[language=TypeScript]
{
  path: 'patient',
  canActivate: [authGuard, roleGuard(['patient'])],
  loadComponent: () =>
    import('./layouts/main-layout/main-layout.component')
      .then((m) => m.MainLayoutComponent),
  children: [
    { path: '', pathMatch: 'full', redirectTo: 'dashboard' },
    {
      path: 'appointments',
      loadComponent: () => import(
        './features/patient/appointments/patient-appointments.component'
      ).then((m) => m.PatientAppointmentsComponent)
    },
    // ... autres routes patient
  ]
}
\end{lstlisting}

\newpage

% ============================================
%  7. PAGES PAR ROLE
% ============================================
\section{Interfaces métier par rôle}

\subsection{Espace Patient}

\begin{table}[H]
\centering
\renewcommand{\arraystretch}{1.4}
\begin{tabularx}{\textwidth}{lX}
\toprule
\rowcolor{primary!10}
\textbf{Page} & \textbf{Fonctionnalités} \\
\midrule
Tableau de bord & 3 tuiles d'accès rapide (RDV, dossier, ordonnances) \\
Rendez-vous & Recherche par spécialité/localisation/date, onglets À venir / Historique \\
Dossier médical & Infos personnelles, vaccinations, historique des consultations \\
Ordonnances & Liste des prescriptions actives, téléchargement PDF, rappels \\
Profil & Édition des informations, sécurité, 2FA \\
\bottomrule
\end{tabularx}
\end{table}

\subsection{Espace Médecin}

\begin{table}[H]
\centering
\renewcommand{\arraystretch}{1.4}
\begin{tabularx}{\textwidth}{lX}
\toprule
\rowcolor{primary!10}
\textbf{Page} & \textbf{Fonctionnalités} \\
\midrule
Tableau de bord & Vue d'ensemble de la journée \\
Planning & Toggle Jour/Semaine/Mois, créneaux colorés par type (consultation, suivi, urgence) \\
Consultations & Liste des patients du jour avec statuts (en attente, en cours, terminée) \\
Patients & Recherche, cartes avec antécédents et dernière visite \\
\bottomrule
\end{tabularx}
\end{table}

\subsection{Espace Secrétariat}

\begin{table}[H]
\centering
\renewcommand{\arraystretch}{1.4}
\begin{tabularx}{\textwidth}{lX}
\toprule
\rowcolor{primary!10}
\textbf{Page} & \textbf{Fonctionnalités} \\
\midrule
Tableau de bord & Accès rapide aux opérations courantes \\
Rendez-vous & Table CRUD complète avec menu d'actions \\
Facturation & KPIs (encaissé, en attente, impayés), génération PDF \\
\bottomrule
\end{tabularx}
\end{table}

\subsection{Espace Administration}

\begin{table}[H]
\centering
\renewcommand{\arraystretch}{1.4}
\begin{tabularx}{\textwidth}{lX}
\toprule
\rowcolor{primary!10}
\textbf{Page} & \textbf{Fonctionnalités} \\
\midrule
Tableau de bord & 4 KPIs (consultations, revenus, occupation, no-show) \\
Établissement & Identité, horaires, spécialités, salles \& équipements \\
Personnel & Gestion des médecins (tarifs, spécialités) et du staff \\
Finance & Graphique barres revenus, tarifs par acte \\
Statistiques & Consultations par praticien, camembert par spécialité \\
\bottomrule
\end{tabularx}
\end{table}

\newpage

% ============================================
%  8. DESIGN SYSTEM
% ============================================
\section{Design System \& UX}

\subsection{Palette de couleurs}

\begin{table}[H]
\centering
\renewcommand{\arraystretch}{1.5}
\begin{tabularx}{\textwidth}{lXl}
\toprule
\rowcolor{primary!10}
\textbf{Couleur} & \textbf{Usage} & \textbf{Code} \\
\midrule
\cellcolor{primary}\color{white}Primary & Actions principales, accents & \texttt{\#0D6EFD} \\
\cellcolor{secondary}\color{white}Secondary & Gradients, accents secondaires & \texttt{\#06B6D4} \\
\cellcolor{accent}\color{white}Accent & Gradient hero, éléments décoratifs & \texttt{\#8B5CF6} \\
\cellcolor{success}\color{white}Success & Confirmations, statuts positifs & \texttt{\#10B981} \\
\cellcolor{warning}\color{white}Warning & Avertissements, en attente & \texttt{\#F59E0B} \\
\cellcolor{danger}\color{white}Danger & Erreurs, annulations & \texttt{\#EF4444} \\
\bottomrule
\end{tabularx}
\caption{Palette principale du design system}
\end{table}

\subsection{Composants partagés}

\begin{itemize}
  \item \textbf{PageHeaderComponent} : entête uniforme (eyebrow + titre + sous-titre + actions).
  \item \textbf{MainLayoutComponent} : shell responsive avec sidebar + toolbar + content.
  \item \textbf{RevealOnScrollDirective} : animation au scroll pour la landing.
\end{itemize}

\subsection{Responsive design}

L'application adopte une approche \textbf{mobile-first} :

\begin{itemize}
  \item Grilles CSS \texttt{auto-fit} / \texttt{auto-fill} avec \texttt{minmax}.
  \item Media queries aux breakpoints 600px, 768px, 900px.
  \item Sidebar transformée en menu burger sur mobile.
  \item Tableaux convertis en cartes verticales sur petits écrans.
\end{itemize}

\subsection{Accessibilité}

\begin{itemize}
  \item Labels et placeholders explicites sur tous les formulaires.
  \item Attributs ARIA sur les boutons icônes (\texttt{aria-label}).
  \item Contrastes conformes WCAG AA.
  \item Navigation au clavier supportée nativement par Angular Material.
\end{itemize}

\newpage

% ============================================
%  9. INTEGRATION BACKEND
% ============================================
\section{Intégration avec le backend}

\subsection{Communication API}

Le frontend dialogue avec une API REST \textbf{Spring Boot 3.5} sur le port
\texttt{8080}, authentifiée par \textbf{JWT (HS512)}.

\begin{lstlisting}[language=TypeScript, caption={Intercepteur JWT}]
export const jwtInterceptor: HttpInterceptorFn = (req, next) => {
  const token = inject(AuthService).getToken();
  if (!token) return next(req);
  return next(
    req.clone({ setHeaders: { Authorization: `Bearer ${token}` } })
  );
};
\end{lstlisting}

\subsection{Endpoints consommés}

\begin{table}[H]
\centering
\renewcommand{\arraystretch}{1.4}
\begin{tabularx}{\textwidth}{>{\ttfamily}lX}
\toprule
\rowcolor{primary!10}
\textbf{Endpoint} & \textbf{Description} \\
\midrule
POST /api/auth/register & Création de compte patient \\
POST /api/auth/login & Connexion et émission de JWT \\
GET /api/users/me & Profil de l'utilisateur connecté (à venir) \\
GET /api/appointments & Liste des rendez-vous (à venir) \\
GET /api/medical-records & Dossiers médicaux (à venir) \\
\bottomrule
\end{tabularx}
\caption{Endpoints principaux}
\end{table}

\subsection{Configuration CORS}

\begin{successbox}[title={Sécurité}]
Le backend autorise uniquement l'origine \texttt{http://localhost:4200} en
développement, avec credentials activés. Cette politique est durcie pour
la production via la variable d'environnement \texttt{environment.apiUrl}.
\end{successbox}

\newpage

% ============================================
%  10. CONCLUSION
% ============================================
\section{Conclusion}

\subsection{Bilan de la partie frontend}

La partie frontend de MediSync illustre l'application des meilleures
pratiques modernes du développement Angular :

\begin{itemize}
  \item Architecture \textbf{Core / Shared / Features} évolutive.
  \item Composants \textbf{standalone} avec \textbf{signaux} pour la réactivité.
  \item \textbf{Lazy-loading} systématique pour des performances optimales.
  \item Design system cohérent avec \textbf{Angular Material 3}.
  \item Authentification robuste avec \textbf{JWT} et \textbf{role-based guards}.
  \item Interface \textbf{responsive} adaptée à tous les écrans.
\end{itemize}

\subsection{Statistiques du projet}

\begin{table}[H]
\centering
\renewcommand{\arraystretch}{1.5}
\begin{tabularx}{\textwidth}{Xr}
\toprule
\rowcolor{primary!10}
\textbf{Indicateur} & \textbf{Valeur} \\
\midrule
Pages métier livrées & 17 \\
Composants standalone créés & 19 \\
Services applicatifs & 1 (AuthService) \\
Guards de protection & 3 (authGuard, roleGuard, guestGuard) \\
Routes lazy-loadées & 17 \\
Lignes de TypeScript & $\approx$ 2 500 \\
Lignes de SCSS & $\approx$ 1 800 \\
\bottomrule
\end{tabularx}
\caption{Volume de code produit}
\end{table}

\subsection{Perspectives}

\begin{itemize}
  \item Intégration de la \textbf{messagerie temps réel} (WebSocket) pour les notifications.
  \item Mise en place d'un \textbf{store global} (NgRx ou Signals Store) si la complexité augmente.
  \item Ajout de \textbf{tests unitaires} et e2e (Jest, Playwright).
  \item Optimisations \textbf{SSR} (Angular Universal) pour le SEO de la landing.
  \item Internationalisation \textbf{i18n} (français, arabe, anglais).
\end{itemize}

\vspace{1cm}

\begin{tcolorbox}[
  enhanced, colback=primary, colframe=primary,
  arc=12pt, drop shadow,
  fontupper=\color{white}\large\bfseries,
  halign=center
]
Merci de votre attention\\[0.2cm]
\normalfont\normalsize Questions \& échanges
\end{tcolorbox}

\end{document}
